\section{Review of the Empirical Literature}

This article aims to unify the literature on residential segregation and neighborhood change in a theoretical model. As such, the model shall be able to explain key stylized facts. A stylized fact is an empirical regularity that abstracts from and lightly theorizes many empirical observations into a general empirical claim especially worthy of explanation \citep{hirschmanStylizedFactsSocial2016}. In this section, I review significant lines of inquiry in urban sociology, demography, and economics to obtain stylized facts about social and economic segregation and its dynamics. However, as most reviewed research considers North America and Europe, the stylized facts and subsequent model will characterize patterns in market democracies. The applicability to other contexts may be possible but would require caution. The stylized facts and the theoretical model will not apply to contexts where housing is not allocated via markets. 


\subsection{Residential Segregation}

Residential segregation refers to the unequal distribution of population characteristics over residential areas beyond random chance \citep{masseyDimensionsResidentialSegregation1988, reardonNewApproachMeasuring2006}. Inequality or heterogeneity within a concept is a necessary condition for segregation. This unequal distribution of population groups in space can have different levels \citep{oecdDividedCitiesUnderstanding2018}, forms \citep{masseyDimensionsResidentialSegregation1988} and geographic scales \citep{leeCensusTractPatterns2008}. Although measurement of segregation is often difficult and relies on arbitrary spatial units \citep{reardonNewApproachMeasuring2006, leeCensusTractPatterns2008, leung-gagneItSurprisinglyDifficult2023}, neighborhoods are usually more homogenous concerning their population compared to entire cities.

Almost every feature of populations is distributed unequally within cities, and this observation holds globally, even though there are considerable differences in levels of segregation between countries and cities \citep{quillianSocioeconomicSegregationLarge2016, oecdDividedCitiesUnderstanding2018, vanhamUrbanSocioEconomicSegregation2021}. Sociologists have given most attention to segregation by race/ethnicity \citep{charlesDynamicsRacialResidential2003}. However, urban populations are also more or less segregated by income and wealth (e.g., \cite{reardonIncomeInequalityIncome2011, oecdDividedCitiesUnderstanding2018, helbigHinterFassadenZur2023}), age and household type and composition (e.g., \cite{owensInequalityChildrenContexts2016, loganBirdsFeatherSocial2016, mutganIncomeInequalityResidential2023}), education and socio-economic status (e.g., \cite{quillianSocioeconomicSegregationLarge2016, helbigHinterFassadenZur2023}), and political opinion and party preference (e.g., \cite{mummoloWhyPartisansNot2017}). 
Relatedly, the type of building and tenure (e.g., \cite{loufPatternsResidentialSegregation2016, oecdDividedCitiesUnderstanding2018, owensBuildingInequalityHousing2019, beaubrun-diantImpactPublicHousing2022}) and housing quality concerning size and amenities of the housing units themselves (for a review, see \cite{chauCriticalReviewLiterature2003}) and of the neighborhood like parks, waterfronts, distance to the center and public transport or disamenities like noise or air pollution are unequally distributed in space (e.g. \cite{leeNaturalAmenitiesNeighbourhood2018, heblichEastSideStoryHistorical2021, ruttenauerEnvironmentalInequalityResidential2021}).

The standard model of the spatial distribution of income has been that of a monocentric city with poor city centers and affluent suburbs \citep{bruecknerLecturesUrbanEconomics2011, rosenthalChangePersistenceEconomic2015}. While this is a good description of most US cities, this model does not hold globally \citep{bruecknerWhyCentralParis1999, quillianSocioeconomicSegregationLarge2016, gaigneWhoLivesWhere2022}. Inequality is a necessary condition for segregation, of course. However, more unequal countries and cities also tend to be more segregated cross-sectionally \citep{quillianSocioeconomicSegregationLarge2016, oecdDividedCitiesUnderstanding2018, vanhamUrbanSocioEconomicSegregation2021} and some scholars argue that income segregation has risen as a result of increased income inequality \citep{reardonIncomeInequalityIncome2011, vanhamUrbanSocioEconomicSegregation2021, mutganIncomeInequalityResidential2023} although this claim is contested \citep{loganUptickIncomeSegregation2018, loganIncomeSegregationWhom2020, molinderLongrunEvidenceStockholm2026}. Segregation is predominantly a phenomenon of the extremes of the income distribution. The richest and, to a lesser extent, the poorest quantiles of the income distributions are most isolated in space \citep{reardonIncomeInequalityIncome2011, quillianSocioeconomicSegregationLarge2016, oecdDividedCitiesUnderstanding2018, mutganIncomeInequalityResidential2023}. Families appear more economically segregated than households without children \citep{owensInequalityChildrenContexts2016, mutganIncomeInequalityResidential2023}.


\subsection{Neighborhood Stability and Change}

Most neighborhoods, especially at the very top and very bottom, maintain their position in the hierarchy of the city over long periods \citep{solariAffluentNeighborhoodPersistence2012, owensNeighborhoodsRiseTypology2012, rosenthalChangePersistenceEconomic2015}. Nevertheless, neighborhoods can change over time, although it is usually a highly localized and rare phenomenon \citep{solariAffluentNeighborhoodPersistence2012, owensNeighborhoodsRiseTypology2012, brown-saracinoExplicatingDividedApproaches2017}. Only a few neighborhoods within a city change at a time, if any at all. 

Gentrification has become the general term for the socio-economic ascent of a neighborhood \citep{owensNeighborhoodsRiseTypology2012}. Initially, gentrification was defined as the ascent of a previously lower-class neighborhood (typically near the city center) with below-average housing quality to a new middle- or upper-middle-class neighborhood with improved housing and infrastructure conditions accompanied by a new dominant lifestyle by an exchange of residents \citep{zukinGentrificationCultureCapital1987}. However, this term has been widened to include more and more diverse cases of neighborhood ascent \citep[530ff]{owensNeighborhoodsRiseTypology2012, brown-saracinoExplicatingDividedApproaches2017}. However, the narrow and wide definitions share that not only do the average incomes and rents in the neighborhood change but that they are associated with socio-cultural changes and changes in the built environment \citep{redfernWhatMakesGentrification2003}. More educated and affluent households with different lifestyles replace the initial residents. The businesses in the area increasingly cater to these new residents. Housing quality and public good provision increase \citep{zukinGentrificationCultureCapital1987}. Gentrification typically occurs in places of disinvestment within areas of at least some prosperity \citep{leyAlternativeExplanationsInnerCity1986, atkinsonEvidenceImpactGentrification2004} or in proximity to already established neighborhoods \citep{guerrieriEndogenousGentrificationHousing2013}.

\begin{figure}
\caption{The Gentrifying Neighbourhood over Time}
\begin{subfigure}{.5\textwidth}
  \includegraphics[width=\linewidth]{/home/malte/Projects/SocEconHousing/images/lit_01_stage_groups.pdf}
  \caption{Group Composition}
\end{subfigure}
\hfill
\begin{subfigure}{.5\textwidth}
  \includegraphics[width=\linewidth]{/home/malte/Projects/SocEconHousing/images/lit_01_stage_capital.pdf}
    \caption{Average Income and Status}
\end{subfigure}

\tiny{Schematic representation of a gentrifying neighborhood over time. (a) is adapted and simplified from \citet[70]{dangschatGentrificationHamburg1991}.}

\label{fig:gentri}
\end{figure}

Gentrification changes a neighborhood economically, socially and in terms of the built environment usually gradually but often fast compared to the usual stability of neighborhoods. The resulting sigmoid curve between two phases of relative stability is typical for neighborhood change \citep[927f]{meenSocialBehaviourBasis2003}. These changes in residents has been described in distinct successions of population groups. Before the beginning of the process of gentrification in a neighbourhood, residents are have low income and low status \citep{zukinGentrificationCultureCapital1987, dangschatGentrificationHamburg1991, blasiusPioneersGentrifiersProcess2016}. In the first stage of gentrification, \textit{pioneers} move in which have low income as well but have a higher socio-economic status. They are often students or creatives \citep{zukinGentrificationCultureCapital1987, blasiusPioneersGentrifiersProcess2016}. They in turn make the neighborhood more attractive, as they change the character and reputation of the neighborhood. The neighborhood then attracts \textit{gentrifiers} with more financial resources which drives up rents and investment. Figure \ref{fig:gentri} provides a schematic overview over this process. Crucially, the social changes of the neighborhood preceed the economic changes.

The literature often relates these successions of residents to forced displacement. However, differentiating between voluntary and forced moves is often complicated, and the amount of residential mobility associated with gentrification is typically low and varies by context \citep{freemanDisplacementSuccessionResidential2005, dingGentrificationResidentialMobility2016, leeGeographyGentrificationResidential2023}. Low-income households show high levels of residential mobility, independent of whether they live in gentrifying neighborhoods. Forced displacement might, therefore, not be necessary for gentrification dynamics; the characteristics of new residents change, but the vacancies are created by normal turnover \citep{freemanDisplacementSuccessionResidential2005}. The potential relationship of gentrification to displacement has spurred a large body of literature on the consequences of gentrification, especially for the displaced but also for the neighborhood and the city as a whole (for a review, see \cite{atkinsonEvidenceImpactGentrification2004}). 

In many ways, neighborhood descent constitutes the mirror image of gentrification, except that decline is a much slower process than ascent, primarily because housing decays at a slow rate \citep{glaeserUrbanDeclineDurable2005}. When a neighborhood declines, it is still attractive to some renters as they find cheap housing with amenities they could not afford elsewhere. As decline progresses, this holds only for continuously poorer households. There is again a replacement cycle where poorer households replace previous residents as infrastructure and housing quality decline. The population structure changes, but the population size does not necessarily decline \citep{guerrieriWithinCityVariationUrban2012}.


\subsection{Housing and Neighborhood Inequality}

Residential segregation and the spatial inequality of housing quality create distinct neighborhoods that affect their resident's lives. Unsurprisingly, the different dimensions of spatial inequalities are related: wealthier households inhabit housing units and neighborhoods of higher quality, whereas disadvantaged households have lower housing and neighborhood quality. Neighborhood inequalities, particularly concerning exposure to pollution and crime, children's school quality, and access to resources and opportunities, cause the reproduction of social inequalities over generations (e.g., \cite{masseyAmericanApartheidSegregation1993, sharkeyIntergenerationalTransmissionContext2008, chettyImpactsNeighborhoodsIntergenerational2018a}).

In recent years, increasing housing costs and stagnating incomes resulted in increasing financial burdens related to housing for low to middle-income households \citep{delucaHousingInsecurityPoor2022}. While more affluent households spend more on housing in absolute terms, poorer households spend a higher proportion of their income on housing. Absolute rent differences between poor and non-poor neighborhoods are small though \citep{desmondPoorPayMore2019}. For many low-income households, housing costs exceed the recommended maximum of 30\% of monthly income. A substantial share of the population even spends more than half their income on housing. In the US, African-Americans and Hispanics are more likely to be housing cost burdened \citep{delucaHousingInsecurityPoor2022}. 

Although they pay a large share of their income on housing, poor households frequently live in housing units with substantial quality issues, such as lacking heating or water, poor building materials, infestations, and mold. Poor households are also more likely to live in overcrowded units. In the US, living in substandard housing is especially prevalent for low-income families and racial minorities. Differences in housing quality between racial groups remain after controlling for income and minority households need more financial resources to attain the same residential outcomes \citep{crowderInterneighborhoodMigrationRace2010, pendallWhyHighPovertyNeighborhoods2016, delucaHousingInsecurityPoor2022}. 

Housing cost burdens and low housing quality are the prime reasons poor households have low housing security and high levels of residential mobility. Landlords (threaten to) evict tenants if they fall behind in rent payments, and tenants are forced to move when units fail to provide necessities. Poor households move often and on short notice. In contrast, affluent households move rarely, and if they do, they move to increase their residential satisfaction or because of life course events such as education, cohabitation, and childbirth \citep{bruchChoiceSetFormation2019, delucaNotJustLateral2020, delucaHousingInsecurityPoor2022}. These forced moves by low-income households often lead them to live in even worse units and neighborhoods, especially when households have been evicted \citep{desmondEvictedPovertyProfit2017, delucaHousingInsecurityPoor2022}. 


\subsection{Summary: Stylized Facts}

\begin{itemize}
\item Residents of a city are not evenly distributed in space. There is always segregation along social, economic, and demographic lines beyond what we expect by random chance.
\item The extremes of the income distribution are the most segregated groups.
\item Housing type and quality and other amenities are also unevenly placed.
\item These different spatial inequalities correlate with each other. Poorer residents typically have lower housing and neighborhood quality.
\item Disadvantaged residents have low housing security and high rates of residential mobility.
\item There is a negative relationship between income and the proportion of income spent on housing.
\item Neighborhoods maintain their rank within the city over long periods. Only a few neighborhoods change their relative position at a time, if at all.
\item Gentrification typically occurs close to established neighborhoods, amenities, or the city center.
\item Economic, socio-cultural, and infrastructural change in the city and neighborhoods are almost always linked. Social change precedes economic change in neighborhood ascent.
\item Gentrification can quickly change a neighborhood, whereas decline is usually slower as built infrastructure creates inertia.
\end{itemize}